\section{Paired cutoffs, exact divisor signs and proved regimes}\label{pca:core}

The comparison in this section concerns a specified part of the signed
coprime-column form for the original unweighted, nonprincipal, two-pivot
Eisenstein rows. We use the physical ideal norms and the original ordered
cutoff. Let \(D,U,H\) be the column, cutoff and capacity parameters, with
\(U^2\le D\le 7U^2/2\) and \(A_1A_2^2\ll H\le D\). The first two row ideals have
dyadic norms \(A_1,A_2\), \(R=A_1A_2\), and the actual primitive conductor of
row \(u\) has norm \(Q_u\asymp R\). The finite bad-prime set \(S\), row type,
and twist are fixed. The same smooth complex profile \(W\) supported on
\([1,2]\) is used on both sides. Ideal labels, including conjugate primes
of the same norm, are retained individually.

\subsection{The selected pairing}\label{pca:pair}
For each good squarefree intermediate M\"obius divisor \(d\), put \(E=Nd\).
Let \(V_d\) be its original row amplitude and \(\mathcal E_d\) the
equal-norm part of the original middle Poisson window. We study selections
of
\begin{equation}\label{pca:cross}
 \frac{2}{D}\operatorname{Re}\sum_d\mu(d)
       \langle \mathcal E_d,V_d-\mathcal E_d\rangle .
\end{equation}
Here the counting inner product is linear in the first argument. This is
one term in the expansion of the original signed norm. In particular the
intermediate \(d\) is different from the gcd of the final column pair.

On an original allocation let \(z=e_0z_0\) with
\(J\le Nz<2J\), let \(e_1\mid\operatorname{rad}(Sd)\) be the other deletion
label, and put \(h=re_0\), with \(r\) the original squarefree short divisor.
Set
\[
 X=\frac{D}{E(Nr)^2},\qquad X'=\frac{X}{(Ne_0)^2},
 \qquad M=\frac{R\,Ne_0\,Ne_1\,J}{X}.
\]
The original middle window has \(X'>B_*\) and \(K_0<Nk\le M D^\eta\),
for its specified \(B_*,K_0,\eta\). The primal test discards the
allocations on which the original radial weight vanishes. An allocation
with \(X'\le B_*\) belongs to the direct component. Its re-expressed
Poisson sum stays in the complement in \eqref{pca:cross}.

On the equality side write \(z_{0,0}=g\ell\) and \(k_0=g\bar\ell\), so
\(Nz_{0,0}=Nk_0\). On the partner side the selected changes in the good
exponents consist of erasing the squarefree cutoff part \(x\), erasing a
part of the good conjugate overlap, and inserting allowed positive
increments at other good primes outside \(d\). A conjugate erasure removes
the full original exponent, which can be two. At an insertion prime the
old residue plus the increment is at most five. Changes at the good
primes of \(d\) keep their exact new residues. For fixed other labels
write \(z_{0,1}=z_*x\) and \(k=Ik_*\). The physical insertion ideal \(I\)
keeps each exponent \(1,\ldots,5\) and is disjoint from the selected
erasure supports and \(d\). Thus \(Nx\asymp B\) and \(NI\asymp K\) are
physical lengths. The two varying raw phases are
\(\psi_u(x)\overline{\psi_u(I)}\), or their conjugates if the pairing is
conjugated. The fixed row factors retain the union of deleted primes,
the literal section, the twist and the primitive \(S\)-values.

The subtraction on the partner imposes \(Nz_*Nx=Nk_*NI\) only in its
original middle window. At fixed \(x\) this fixes at most one integer
norm for \(I\), and the reverse count fixes at most one norm for \(x\).
An integer norm supports at most divisor-many ideals on the polynomial
clip. With the physical inverse \(1/Nx\), the paired count therefore
costs \(\min(1,K/B)\) up to subpowers. The two half radial powers leave
\[
 (B/K)^{1/2}\min(1,K/B)=\sqrt{\min(B/K,K/B)} .
\]
The remaining harmonic weights give \(O(R/E)\) for this subtraction at
each \(d\), with the inherited profile and subpower costs. Both original
endpoints remain in the count. Equal norms in the direct and other
components remain in the pairing.

\subsection{Target bounds on the selected cross}\label{pca:bounds}
The bounds have the common cost
\(O_{\epsilon,S,\nu,u_0}((DH)^\epsilon\mathcal M_{24}(W)^2)\), where
\(\mathcal M_{24}\) controls the indicated smooth profile derivatives.
The live partner primal test supplies \(E\ll D/B\) and partner radial
scale \(y'_1\gg Nk/R\gg K/R\). The proof first sums each actual column
array. It then uses the row comparisons and the original label sums.

At the mixed scales \(A_1\asymp D^{5/12},\, A_2\asymp D^{1/6}\), the
selected normalized cross has upper power \(57/80\) over all physical
ratios and all \(d\). If \(H\ll D^{3/4}\), a perturbed primitive sharp
moment gives the target power \(3/4\) for the same selected cross on
\[
 \frac7{24}\le\log_D A_1\le\frac{439}{912},\qquad
 A_1A_2^2\ll H.
\]
The second pivot need only satisfy the stated original admissibility
and capacity. The conductor comparison uses the actual \(R\).

At the pure critical scales \(A_1\asymp R\asymp D^{3/4}\), \(A_2\asymp1\),
the target power \(3/4\) holds over all \(d\) in each of the following
physical sectors:
\begin{itemize}
\item \(K\gtrsim D^{1089/1420}\), with any erasure length \(B\).
\item \(B\gtrsim D^{1355/1482}\), with any insertion length \(K\).
\item \(B\gtrsim R\) and \(K\lesssim D^{921/1652}\).
\end{itemize}
The flexible central windows give further short and long threshold
curves. There are also separate comparisons for a small fixed residual
label product and for a small \(d\)-prefix. The comparable pure
all-divisor comparison has power \(55/68\), and the target holds on
\(Nd\le D^{9/68}\) in that sector. Each assertion concerns the selected
changes described above. The remaining pure rectangles and the
coherently averaged signed divisor sum are open.

\subsection{Ordered assignment signs}\label{pca:sign}
Fix an ordered factorization \(\zeta=\zeta_1\zeta_2\) of the actual
squarefree cutoff label \(\zeta=e_0z_0\) on either side, with
\((d,\zeta)=1\) and \(N\zeta_i\le U\). Write \(T_i=U/N\zeta_i\). The
assignment sum multiplying \(\mu(\zeta)\) is
\[
 F_d=\sum_{\substack{a,b\mid d,\ (a,b)=1\\Na\le T_1,\ Nb\le T_2}}\mu(a)\mu(b).
\]
Complete the \(b\)-sum and then subtract its excluded terms. With
\(m=d/b\) in the subtraction, squarefreeness gives the exact identity
\begin{equation}\label{pca:ordered}
 \frac{F_d}{\mu(d)}=
 \mathbf1_{E\le T_1}+\mathbf1_{E\le T_2}-1
 -\sum_{\substack{m\mid d\\T_1<Nm<E/T_2}}
      \sum_{\substack{a\mid m\\Na\le T_1}}\mu(m/a).
\end{equation}
The complete \(b\)-sum forces \(a=d\). The unit \(m=1\) gives the
indicator at \(T_2\), while the nonunit terms with \(Nm\le T_1\) vanish by
the full inner M\"obius sum. Both inequalities of the remaining band
are strict. Below the product cap the band is empty. The three possible
assignment values are \(\mu(d),0,-\mu(d)\) according as both, just one,
or neither individual capacity reaches \(E\).

Fix the row, both ordered factorizations and the primal and dual
labels. Select both \(e_1\) labels to be fixed \(S\)-supported ideals,
take the good changed-residue ideal on \(d\) to be the unit, and require
\(d\) to avoid the full fixed prime union of the row and the primal and
dual labels. When \(E N\zeta\le U^2\) on both sides, their capacity step
functions multiply the original outer M\"obius sign. The fixed-field
deleted M\"obius estimate supplies subpower cancellation for the
resulting fixed-label scalar on \(Nd\asymp Y\) with \(Y\ge D^\kappa\) for
a fixed \(\kappa>0\). In the pure comparable block, a physical \(Nx\)-interval of
width \(O(BD^{-1/8})\) gives the target after the full divisor sum by a
separate sparse erased-column estimate. This includes the narrow
interval immediately above the product cap on the actual partner
allocation. The wide correction and the remaining coherent signed
estimate require additional arguments. The whole signed and unsigned
powers \(13/12\) and \(9/8\), the independent determinant comparison,
and the scope-specific zero-location obligations remain as in the
preceding sections.
